\section{Level 1.2b: Ecologically Stratified Rediscovery}
\label{sec:level12b}

\subsection{Exact accounting and ledger validation}
\label{sec:level12b-ledger}

Before any feasibility ratio could be trusted, Level 1.2b required its own
ecological ledger --- a full supply/demand/net-zero-transfer
reconciliation of total system resources (every patch's stock, every
agent's resources, and every active rule's fund) --- to close to zero
within a documented floating-point tolerance
(\cref{sec:appendix-accounting} gives the full identity and category
list). The ledger was first validated only at short scale (population 16,
3 patches, 40 steps) and marked complete on that basis. While preparing
the milestone's actual production-scale design (population 32, 4 patches,
700 steps), the identical configuration preset showed a residual of
$9.14$, not approximately zero --- the ledger-closure gate blocked the
main discovery run pending investigation, exactly as its own preregistered
stop condition required.

All six root causes shared one property: none could manifest at short
scale. They were, in order of discovery: (1) a sign-flip in every
capped-resource-deduction site, so that a target agent already in a
negative-resource distress state would have resources incorrectly
\textit{restored} rather than further deducted; (2) a monitoring-cost
event stamped onto every violation observed in a round rather than once
per round, double- and triple-counting a cost paid once; (3) a per-active-
link coordination maintenance fee charged every step but previously logged
only when a separately gated activity roll fired, silently under-recording
the fee on steps where it did not; (4) a rule-identifier/patch-identifier
key-space conflation in dissolution routing that could silently corrupt
institutional state via unrelated numeric coincidences; (5) an amendment-
triggered dissolution path that never removed the dissolved rule from the
live registry, leaving a stale entry whose fund kept counting toward
system totals; and (6) a zero-recipient fund-disposition case not treated
as a loss. After all six fixes, the production-scale configuration
reconciles to a residual on the order of $10^{-10}$ for every seed tried
(\cref{fig:l12b-ledger}); across the milestone's full 576-run discovery
design, the maximum observed residual was $5.9\times10^{-9}$ and the mean
$7.1\times10^{-10}$, confirming the fix generalizes across the full design
space rather than only the configuration that first exposed it. These were
genuine accounting bugs, invisible below production scale, found and
fixed by this milestone's own mandatory preflight audit --- not
generative-behavior changes: every fix is either a bookkeeping correction
or a corrected cleanup of already-dissolved state, confirmed
behavior-preserving against the full pre-existing test suite at every step
of the fix sequence.

\subsection{Ex-ante feasibility measures}
\label{sec:level12b-feasibility}

Two ex-ante ecological feasibility ratios were defined and validated
before being used to stratify the main design. The
\textit{structural feasibility ratio} is a pure function of configuration
(no simulation, no randomness): maximum-sustainable-yield regeneration
(patch regeneration rate times capacity over four, the logistic-growth
maximum at half capacity, summed across patches) divided by population
size times the bare mandatory subsistence cost --- deliberately excluding
every discretionary or institutionally contingent cost, so it never
assumes optimistic endogenous cooperation in either direction. The
\textit{expected-pressure ratio} uses the same numerator but adds expected
behavioral extraction demand, evaluated at configured mean trait values
and a preregistered neutral occupancy of one-half. A feasibility-map
experiment (120 runs: a nine-column Latin Hypercube design over ecological
parameters, at least two population/patch-count combinations, institutional
configuration held fixed) validated both ratios against realized ecological
balance. The expected-pressure ratio predicted realized ecological pressure
substantially better than the structural ratio (Pearson $r\approx0.84$
versus $r\approx0.48$; \cref{fig:l12b-feasibility-ratios}), an expected,
un-tuned result given that only the expected-pressure ratio incorporates
behavioral extraction demand. The feasibility map itself returned a
genuinely informative, non-degenerate three-way outcome split (50 other,
52 collapse, 18 saturation out of 120 runs; \cref{fig:l12b-feasibility-strata}).
Tertile cut points on the expected-pressure ratio, computed from this
design's own ecological draws and frozen before any discovery run, define
the three ecological strata (low, medium, high) used to stratify the main
joint design.

\subsection{Stratified experimental design}
\label{sec:level12b-design}

The main discovery design (Work Package E) crosses a freshly sampled
ecological feasibility space with a freshly sampled 21-dimensional
institutional Latin Hypercube design (60 institutional points, verified
free of correlated parameter "packages") and two initial-condition
families --- an institution-free default state, and a heterogeneous
proto-institution injected independently per patch --- stratified only by
the ex-ante expected-pressure ratio, never by any realized simulation
outcome. The preferred design (48 joint points $\times$ 6 seeds
$\times$ 2 initial-condition families = 576 runs) was executed in full: 576
of 576 runs completed with zero failures, stratum spread 17/16/15 points
(204/192/180 runs) across low/medium/high, both initial-condition families
balanced exactly at 288/288, and every run's ledger residual under the
$10^{-4}$ tolerance (\cref{tab:l12-vs-l12b-design} compares this design
directly against Level 1.2's own 24-run design).

\subsection{Discovery and confirmation samples}
\label{sec:level12b-discovery-confirmation}

The main discovery pipeline (Work Package J) reused Level 1.2's own
failure/viability classification unmodified and added a five-label
failure-trajectory classification. On the full 576-run pool: all-outcome
distribution collapse 326, viable 161, fragmented 89, every other label
zero (\cref{fig:l12b-outcome-distribution}); failure-trajectory
distribution sustained collapse 305, no collapse observed 194, transient
dip recovered 56, late-onset collapse 21 --- most collapse outcomes were
sustained across a run's duration, not a last-instant artifact. All 576
runs cleared the non-normative sustained-system data-sufficiency screen,
comfortably above both the preferred (60) and hard-minimum (40) gates.

Clustering and structure discovery (Work Package K) reused Level 1.2's own
primitives unmodified on this much larger pool and reached headline
$k=4$ with verdict \texttt{continuum} (three $k$'s within 0.05 silhouette
of the best; \cref{fig:l12b-discovery-confirmation}). Supporting
diagnostics: bootstrap stability mean adjusted Rand index 0.598 ($k=4$,
30 resamples); mean pairwise algorithm-agreement adjusted Rand index
0.192 (dragged down by \texttt{HDBSCAN} finding zero clusters and 576
noise points on this feature space, while the Gaussian-mixture/
hierarchical pairwise agreement alone was 0.576); missingness-strategy
robustness 0.591 (\texttt{complete\_case}) and 0.653
(\texttt{missingness\_indicator}); cluster assignment essentially
unrelated to initial-condition family (adjusted Rand index 0.0008,
ruling out "clustering is just rediscovering which initial condition a
run started from"); temporal-window stability 0.141 (end-of-run versus
mid-run cluster assignments diverge substantially --- cluster membership
is not stable across a single run's own duration, a genuine finding in
its own right); parameter-neighborhood stability 0.626; and cluster
assignment essentially unrelated to the failure/success split (adjusted
Rand index 0.0, ruling out "clustering is just rediscovering collapse
versus not").

Per the preregistered null-replication rule, a \texttt{continuum} verdict
is treated exactly as a no-candidate finding: no discrete institutional
class was named, bounded, or given an expected direction anywhere in the
confirmation design. What was tested instead was the one concrete object
discovery did produce --- the headline four-cluster partition --- stress-
tested for out-of-sample recurrence on a fresh, disjoint 384-run
confirmation pool (24 joint points $\times$ 8 seeds $\times$ 2
initial-condition families, a disjoint fresh-seed base). The confirmation
pool's own all-outcome distribution (collapse 245, viable 86, fragmented
53) closely tracked the discovery pool's; its own independent clustering
refit reached headline $k=2$, not 4, with verdict \texttt{unstable} (mean
algorithm-agreement adjusted Rand index 0.276, bootstrap mean adjusted
Rand index 0.182, both below the 0.3 stability bar) --- never
\texttt{multi\_config}, so the primary preregistered null expectation
held. Assignment stability between the frozen discovery fit and the
confirmation pool's own independent labels was moderate (adjusted Rand
index 0.527), and feature-profile stability across matched clusters was
high (mean cosine similarity 0.948) --- both consistent with a continuum
rather than either pure noise or a sharply discrete structure. The overall
confirmation verdict was \texttt{null\_replicated}.

\subsection{Cluster stability and held-out replication}
\label{sec:level12b-stability}

\cref{fig:l12b-stability-metrics} summarizes the discovery pool's own
stability diagnostics; \cref{fig:l12b-discovery-confirmation} plots the
discovery and confirmation silhouette curves side by side, visually
showing both the discovery pool's shallow, multi-peaked profile and the
confirmation pool's independent, differently shaped profile --- the
headline $k$ itself is unstable across independent samples, which is
direct evidence for reading the discovery pool's own \texttt{continuum}
verdict as gradedness rather than an artifact of one sample.

A genuine seed-clustered bootstrap (resampling whole ecological-and-
institutional joint-point groups, correctly accounting for the multiple
seeds nested within each point, rather than treating every seed-row as an
independent draw) reproduced the discovery pool's own stability estimate
closely (seed-clustered mean adjusted Rand index 0.591 versus the flat
bootstrap's 0.598), confirming the original result was not an artifact of
a less conservative resampling unit. Bootstrap 95\% confidence intervals
were computed for the headline point estimates: collapse fraction 0.566
[0.526, 0.608] ($n=576$); population-fraction recovery 0.0 [0.0, 0.0]
($n=32$, a tight zero-width interval, itself informative); ecological-
fraction recovery 0.875 [0.75, 0.97]. All four pooled run sets (Work
Packages E, L, M, N) were verified pairwise seed-disjoint.

\subsection{Scale tests and causal ablations}
\label{sec:level12b-scale-causal}

A 32-run scaling-robustness grid (2 population sizes $\times$ 2 patch
counts $\times$ 2 network-modularity levels $\times$ 4 seeds, scale
invariance explicitly not required) found collapse remained dominant at
smaller, structurally varied scales (25 collapse, 4 fragmented, 3 viable;
\cref{tab:outcome-distribution}) and exposed a genuine, disclosed
dissociation between two recovery dimensions after a standardized shock:
population-fraction recovery was 0 of 32 (0\%) at every scale tested,
while ecological-fraction recovery was 28 of 32 (87.5\%). This
dissociation is precisely why no "attractor" label was assigned anywhere
in this milestone: an attractor claim would require majority recovery on
\textit{both} tracked dimensions, and population's own 0\% alone rules it
out (\cref{fig:l12b-causal-ablation} places this alongside the causal
decomposition below; see \cref{sec:appendix-tables} for the full
scaling-robustness table).

A thirteen-factor, 208-run matched-seed causal decomposition (eight
matched-seed pairs per factor, \texttt{rng\_mode="mechanism\_%
substreams"}) moved beyond Level 1.2's own two-factor ablation to test
every major mechanism family's causal contribution to ecological outcome
directly. The two largest ecological-fraction effects were
\texttt{subsistence} (mean paired difference $+0.198$; collapse
probability fell from 1.00 to 0.62 as subsistence cost rose from 0.1 to
0.6 --- counterintuitive, reported as found, not reinterpreted) and
\texttt{migration\_on\_off} (mean paired difference $-0.110$; migration on
worsens the ecological outcome, collapse probability rising from 0.38 to
0.88, directly replicating Level 1.2's own Section 13 migration finding
at roughly 24 times the scale). \texttt{migration\_congestion} and
\texttt{patch\_regen} each contributed a further $\approx+0.11$ and
$+0.10$ respectively. The smallest effect, \texttt{transmission\_cost},
showed essentially no causal signal (mean paired difference $-0.002$,
zero outcome-label transitions in 8 pairs) --- the institutional-
transmission-probability lever this milestone tested does not move
collapse risk in this design. \cref{tab:causal-ablation} reports the full
ranked table; results were never fed back into any already-frozen
artifact, honoring the milestone's own recalibration-prohibition stop
condition.

\subsection{The final continuum/null verdict}
\label{sec:level12b-verdict}

Work Package O's synthesis, applying a closed seven-label cautious
vocabulary and a fixed rule that no "attractor" label may be assigned
without direct evidence of majority recovery on every tracked dimension,
classified the milestone's overall result as \texttt{continuum}, inherited
directly from Work Package K's own discovery verdict, and assigned no
attractor label anywhere. Of six candidate explanations considered for why
cross-run phenotype similarity is observed at all, the evidence was mixed
by design rather than uniform: consistent with temporary (within-run)
similarity and with independent structural convergence (a proxy measure,
not a literal lineage-based test); consistent with environmentally forced
similarity for the three substantial clusters; against measurement
compression (five of eleven raw principal components were needed to reach
90\% explained variance) and against failure-state similarity (adjusted
Rand index against a collapse/non-collapse proxy of 0.0); and
\texttt{shared\_lineage\_persistence} was not evaluated at all, disclosed
explicitly as an evidence gap rather than silently omitted, since
evaluating it would require a fourth full 576-run re-simulation with
per-run lineage identifiers exported, which this milestone judged
low-value given that the continuum and null-replication findings already
independently established that no discrete, well-separated candidate
structure exists for shared ancestry to explain in the first place.
